\phantomsection
\addcontentsline{toc}{section}{Lời cảm ơn}
\section*{Lời cảm ơn}

Lời đầu tiên, chúng em xin trân trọng gửi lời cảm ơn sâu sắc nhất đến thầy Mai Tiến Dũng, giảng viên phụ trách môn học Thị giác máy tính (CS231), đã tận tình truyền đạt những kiến thức chuyên môn quý báu, định hướng phương pháp nghiên cứu và luôn hỗ trợ, giải đáp thắc mắc cho chúng em trong suốt quá trình học tập và thực hiện đề tài.

Chúng em cũng xin chân thành cảm ơn Ban Giám hiệu Trường Đại học Công nghệ Thông tin, Đại học Quốc gia TP. Hồ Chí Minh cùng quý thầy cô Khoa Khoa học Máy tính đã tạo điều kiện học tập tốt nhất, cung cấp môi trường học thuật chuyên nghiệp để chúng em hoàn thành tốt đồ án môn học này.

Mặc dù nhóm chúng em đã nỗ lực hết mình với tinh thần làm việc nghiêm túc và cầu thị, song do giới hạn về mặt thời gian và kiến thức thực tế, báo cáo đồ án khó tránh khỏi những thiếu sót nhất định. Nhóm chúng em rất mong nhận được những nhận xét, góp ý quý báu từ thầy để đề tài ngày càng được hoàn thiện và có tính ứng dụng thực tiễn cao hơn.

\vspace{0.6cm}
\begin{flushright}
\textit{TP. Hồ Chí Minh, Học kỳ 1 -- Năm học 2026 -- 2027}\\[0.2cm]
\textbf{Nhóm sinh vi\^en thực hiện}\\[0.15cm]
Từ Lê Tuấn Anh (24520133)\\
Nguyễn Thanh Tân (24521581)
\end{flushright}

\vspace{0.8cm}
\subsection*{Bảng phân công nhiệm vụ và trách nhiệm thực hiện}
\addcontentsline{toc}{subsection}{Bảng phân công nhiệm vụ và trách nhiệm thực hiện}

Nhằm đảm bảo tính minh bạch, độc lập và chuyên sâu trong nghiên cứu, đồ án được phân chia nhiệm vụ đồng đều giữa hai thành viên ($50\% - 50\%$) theo hai trục tiếp cận bổ trợ cho nhau:

\begin{table}[H]
\centering
\small
\caption{Phân công nhiệm vụ chi tiết và tỷ lệ đóng góp của các thành viên.}
\label{tab:phan_cong}
\begin{tabularx}{\textwidth}{|p{3.2cm}|p{3.5cm}|Y|}
\hline
\textbf{Thành viên} & \textbf{Mô-đun phụ trách} & \textbf{Nội dung kỹ thuật cụ thể} \\
\hline
\textbf{Từ Lê Tuấn Anh}\newline MSSV: 24520133 & Mô-đun Định vị (YOLO11s) \& Thị giác truyền thống & 
-- Tiền huấn luyện (CCTSDB) và tinh chỉnh mô hình YOLO11s 1-Class (\texttt{sign}) để cắt ảnh ROI.\newline
-- Thiết kế chuỗi đặc trưng đa tầng $S0 \rightarrow S3$ (HSV Circular Shift, Tỷ lệ khung hình Aspect Ratio, Lưới không gian $2 \times 2$, HOG 9 hướng gradient).\newline
-- Tối ưu hóa bộ phân loại Support Vector Machine (SVM) với kernel RBF. \\
\hline
\textbf{Nguyễn Thanh Tân}\newline MSSV: 24521581 & Mô-đun Học sâu (CNN) \& Khung Đối chuẩn Đa tầng & 
-- Nghiên cứu, thiết kế và tối ưu mạng nơ-ron tích chập Enhanced LeNet-5 (PyTorch).\newline
-- Thiết lập pipeline tăng cường dữ liệu (Data Augmentation), tối ưu hóa hàm mất mát Cross-Entropy và Adam Optimizer kết hợp StepLR Scheduler.\newline
-- Xây dựng cơ chế chia tách Block Split theo chuỗi video liên tục (Zero Data Leakage), đo lường đối chuẩn 2 tầng (ROI \& Scene Benchmark). \\
\hline
\end{tabularx}
\end{table}

\textit{Ghi chú:} Trong toàn bộ nội dung báo cáo sau đây, các phương pháp và kết quả được trình bày dưới góc nhìn khoa học khách quan và thống nhất của toàn bộ nhóm nghiên cứu.
